% !TEX root = ../main.tex
\chapter{Criando processadores: o Hub de Processadores}
\label{cap:criando-processadores}

Com o projeto aberto, o botão \botao{Hub de Processadores} na toolbar abre o formulário que dá vida a um novo processador SAPHO. É aqui que as escolhas de arquitetura do Capítulo~\ref{cap:processador-sapho} viram campos concretos.

\figurapendente{Modal Criar Hub de Processadores}{O formulário completo, preenchido para o processador \code{media\_movel} do exemplo condutor (Total de Bits 16, Mantissa 10, Expoente 5).}

\section{Os campos do formulário}

\begin{longtable}{@{}p{0.30\linewidth}llp{0.34\linewidth}@{}}
\toprule
\textbf{Campo} & \textbf{Diretiva} & \textbf{Padrão} & \textbf{Validação} \\
\midrule
\endhead
Nome do Processador & \code{\#PRNAME} & \code{procTest\_00} & Letras, números, \code{\_} e \code{-}; sem espaços ou acentos. \\
\multicolumn{4}{@{}l}{\emph{Seção Formato Numérico}} \\
Total de Bits & \code{\#NUBITS} & 23 & Inteiro $>0$ e igual a Mantissa + Expoente + 1. \\
Ganho & \code{\#NUGAIN} & 128 & Inteiro $>0$ e \textbf{potência de 2}. \\
Bits da Mantissa & \code{\#NBMANT} & 16 & Inteiro $>0$. \\
Bits do Expoente & \code{\#NBEXPO} & 6 & Inteiro $>0$. \\
\multicolumn{4}{@{}l}{\emph{Seção Pilhas e Portas}} \\
Tamanho da Pilha de Instruções & \code{\#SDEPTH} & 5 & Inteiro $>0$. \\
Tamanho da Pilha de Dados & \code{\#NDSTAC} & 5 & Inteiro $>0$. \\
Portas de Entrada & \code{\#NUIOIN} & 1 & Inteiro $>0$. \\
Portas de Saída & \code{\#NUIOOU} & 1 & Inteiro $>0$. \\
\bottomrule
\end{longtable}

O formulário valida enquanto você digita: um campo inválido ganha borda vermelha e um tooltip explicando a regra (\enquote{Total de Bits deve ser igual a Bits da Mantissa + Bits do Expoente + 1}, \enquote{Ganho deve ser potência de 2}). O botão \botao{Gerar Processador} só habilita quando tudo está válido. Note que o padrão de fábrica já é consistente: $23 = 16 + 6 + 1$.

Como escolher os valores:
\begin{itemize}
  \item \textbf{Total de Bits} dimensiona os inteiros e o consumo de hardware; 16 bits bastam para muitos sinais de instrumentação, 23 (padrão) dá folga com \emph{float} de boa precisão.
  \item \textbf{Mantissa/Expoente} definem a precisão e a faixa do ponto flutuante (Seção~\ref{sec:float-sapho}); se o seu programa só usa inteiros, os valores têm pouco efeito prático, mas a igualdade estrutural continua obrigatória.
  \item \textbf{Pilha de Instruções} limita a profundidade de chamadas de função aninhadas; \textbf{Pilha de Dados} limita a complexidade das expressões. Os padrões atendem programas típicos.
  \item \textbf{Portas} de entrada/saída: uma por fluxo de dados que o processador troca com o mundo (Capítulo~\ref{cap:cmm-io-stdlib}).
  \item \textbf{Ganho}: o divisor fixo da função \code{norm()}; relevante apenas se você a usa.
\end{itemize}

\section{O que é gerado}

Ao clicar em \botao{Gerar Processador}, a AURORA cria dentro do projeto:

\begin{lstlisting}
<projeto>/<NomeDoProcessador>/
|-- Software/<NomeDoProcessador>.cmm
|-- Hardware/        (vazia; recebera o Verilog e as memorias)
`-- Simulation/      (vazia; recebera o testbench)
\end{lstlisting}

O \file{.cmm} inicial vem com o cabeçalho de diretivas preenchido com os valores do formulário e um \code{main()} vazio à sua espera:

\begin{lstlisting}[style=cmm]
#PRNAME media_movel
#NUBITS 16
#NDSTAC 5
#SDEPTH 5
#NUIOIN 1
#NUIOOU 1
#NBMANT 10
#NBEXPO 5
#NUGAIN 128

void main()
{
    // Øk. Você criou um processador em C±, mas e agora?
}
\end{lstlisting}

E agora, o Capítulo~\ref{cap:cmm-fundamentos}: escrever o programa. O processador também é registrado no \file{.spf} (nomes duplicados são bloqueados, sem distinção de maiúsculas).

\section{Configurações de simulação do processador}
\label{sec:proc-config}

Ao lado do botão de compilar há a engrenagem \botao{Configurações de simulação do processador}, que abre um popover com três ajustes \textbf{por processador}:

\begin{longtable}{@{}lp{0.52\linewidth}l@{}}
\toprule
\textbf{Campo} & \textbf{Efeito} & \textbf{Padrão} \\
\midrule
\endhead
\textbf{Clock frequency (MHz)} & Frequência do clock do \emph{testbench} gerado. & 100 \\
\textbf{Number of clocks} & Quantos ciclos a simulação executa antes de encerrar. & 2000 \\
\textbf{Show arrays in GTKWave} & Emite cada elemento de array como um sinal visível nas ondas (equivale à flag \code{--array} do compilador). & desligado \\
\bottomrule
\end{longtable}

O popover também exibe o \textbf{tempo estimado de simulação} (ciclos divididos pelo clock, em $\mu$s). Se a sua simulação termina antes de o programa produzir os resultados, aumente o número de ciclos aqui.

\section{Renomear e excluir processadores}

Pelo menu de contexto do processador na árvore. Renomear é uma operação completa: a pasta é movida, os artefatos (\file{.cmm}, \file{.asm}, \file{Hardware/<nome>.v}, \file{Simulation/<nome>_tb.v}) são renomeados e a diretiva \code{\#PRNAME} dentro do fonte é corrigida --- tudo numa ação. As mesmas regras de nome do formulário se aplicam. Excluir remove a pasta do processador e o registro no \file{.spf}.

\begin{aviso}
O nome do processador e a diretiva \code{\#PRNAME} precisam concordar --- é assim que a toolchain conecta o fonte aos artefatos. Renomeie sempre pela AURORA, que mantém os dois lados em sincronia.
\end{aviso}
